% Candidate text only. This file does not modify main.tex.

\subsection{Robustness of the finite continuation}

The local differential structure does not uniquely determine a finite
coordinate.  We therefore embedded equation~(14) in the matched-tangent
power family
\begin{equation}
 g_r(x;x_0)=
 \begin{cases}
 [(x/x_0)^r-1]/r, & r\ne0,\\
 \ln(x/x_0), & r=0,
 \end{cases}
 \qquad
 q_{r,s}=\ln\frac{C_n}{C_n^0}
 +A g_r(V_{\rm cem};V_0)
 +\Gamma g_s(\phi_b;\phi_b^0).
\end{equation}
Every member has the same value and the same first derivative
$A/V_0-\Gamma/\phi_b^0$ at the reference point.  The alternatives were
evaluated with the same $A$ and $\Gamma$, observations, weights and ridge
points; no coefficient was refitted to the held-out trajectory.  We also
integrated the local profile slope
\begin{equation}
 \frac{d\ln C_n}{dV_{\rm cem}}
 =-\frac{s_{\ell}^{\mathsf T}P s_V}
         {s_{\ell}^{\mathsf T}P s_{\ell}},
\end{equation}
recomputing the metric along the path.  This numerical curve is a
differential reference rather than a same-complexity two-coefficient model.

\begin{table}
\centering
\caption{Matched-tangent finite-coordinate audit.  Pooled drift is the
maximum absolute coordinate-ratio error on the complete raw structural
ridge.  Level error is the worst absolute leave-one-trajectory-out error
over both directions and raw/adjusted metrics.  Shape RMS is averaged over
the four held-out cases between the two operating points.}
\begin{tabular}{lrrr}
\toprule
Continuation & Pooled drift (\%) & LOTO level (\%) & Shape RMS in $\ln C_n$\\
\midrule
log contact + log HS (selected) & 1.80 & 3.29 & 0.00948\\
log contact + linear HS         & 2.17 & 3.54 & 0.0112\\
log contact + reciprocal HS     & 1.74 & 3.03 & 0.00776\\
linear contact + log HS         & 97.4 & 541  & 0.748\\
first-order local tangent       & 99.9 & 543  & 0.750\\
integrated local projection     & 0.0011 & 2.43 & 0.00210\\
\bottomrule
\end{tabular}
\end{table}

The result supports a strongly curved, approximately logarithmic contact
continuation and rejects extension of the local tangent over the observed
parameter separation.  It does not uniquely select the finite HS factor:
linear and reciprocal HS continuations perform comparably because the HS
term is smaller and $\phi_b$ varies over a restricted relative range.  We
therefore interpret equation~(14) as a stable, low-order physics-factored
coordinate, not as the unique finite state compatible with the local
Jacobian.

\section{Analytic--numerical self-consistency of the local slope}

For either endpoint modulus, write
$m=\partial\ln S/\partial\ln a_c$.  Scheme~1 gives
\begin{equation}
 d\ln K_b=\left(1-\frac{m_K}{4}\right)d\ln C_n
          +\frac{m_K}{4}d\ln V_{\rm cem},
 \qquad
 A_K=\frac{m_K}{4-m_K},
 \qquad
 \beta_K=\frac{A_K}{V_0}.
\end{equation}
The shear endpoint has the same form with the effective exponent
$m_G=\partial\ln G_b/\partial\ln a_c$, including its normal and tangential
terms.  At the pooled point, the sample-median exponents are
$m_K=0.940869$ and $m_G=0.925754$, which give
$\beta_K=24.515$ and $\beta_G=24.002$.  Within Scheme~1, these endpoint
values provide a contact-only benchmark for the observable coordinate;
they are neither an independent prediction of its complete slope nor a
validation against external contact measurements.

To propagate the result to the data, let $u_a$, $u_C$ and $u_b$ denote the
derivatives of the stacked elastic observations with respect to
$\ln a_c$, $\ln C_n$ at fixed $a_c$, and $\phi_b$, respectively, after
division by the data standard deviations.  The target column for
$\ln C_n$ is
\begin{equation}
 j_\ell=u_C-\frac14u_a.
\end{equation}
With $P=I$ for the raw metric and
$P=I-J_\eta(J_\eta^{\mathsf T}J_\eta+I)^{-1}J_\eta^{\mathsf T}$ after
nuisance adjustment, the two finite-coordinate coefficients are
\begin{align}
 A &=
 \frac{j_\ell^{\mathsf T}P(u_a/4)}
      {j_\ell^{\mathsf T}Pj_\ell},\\
 \Gamma &=
 \phi_b\frac{j_\ell^{\mathsf T}Pu_b}
               {j_\ell^{\mathsf T}Pj_\ell},\\
 \beta &= \frac{A}{V_0}-\frac{\Gamma}{\phi_b}.
\end{align}
The raw observable-weighted contact term is 24.120 and the moving-endpoint
HS contribution is $-3.849$, giving $\beta=20.272$.  Nuisance projection
changes the two terms to 24.138 and $-4.176$, giving $\beta_{\rm adj}=19.962$.
Thus, within the same forward model, observable metric and local projection,
the analytic endpoint decomposition and the directly calculated ridge slope
are numerically self-consistent once the HS porosity-path contribution is
included.

Applying the same calculation to all 400 primary 20-m moving-block samples
recovers the stored coefficients because both routes use the same projected
local derivatives.  This is an analytic--numerical implementation check, not
an independent physical prediction.  At the fixed pooled $V_0$, the adjusted
total has median 20.063 and conditional 2.5--97.5 percentile range
$[19.136,20.517]$; $A_{\rm adj}$ has median 0.3041 and range
$[0.2901,0.3114]$, while $\Gamma_{\rm adj}$ has median 1.6207 and range
$[1.5528,1.6680]$.  With only about six block-length units per trajectory,
these ranges are conditional diagnostics with uncertain finite-sample
coverage.  The endpoint exponents constrain the contact contribution
underlying $A$; $\Gamma$ arises from the moving-porosity/HS derivative.
